\documentclass[11pt]{article}
\usepackage[margin=1in]{geometry}
\usepackage{enumitem}
\usepackage{longtable}
\usepackage{booktabs}
\usepackage{amsmath,amssymb}
\usepackage{tcolorbox}
\tcbuselibrary{breakable}
\usepackage{hyperref}
\hypersetup{colorlinks=true, linkcolor=blue, urlcolor=blue}

\setlist[itemize]{itemsep=2pt, topsep=3pt}
\setlist[enumerate]{itemsep=6pt, topsep=3pt}
\setlength{\parindent}{0pt}
\setlength{\parskip}{5pt}

\newtcolorbox{defbox}{breakable, colback=blue!4, colframe=blue!55!black,
  title={Key Definitions}, fonttitle=\bfseries}
\newtcolorbox{bridgebox}{breakable, colback=green!5, colframe=green!45!black,
  title={Connecting to Prior Learning}, fonttitle=\bfseries}
\newtcolorbox{notebox}{breakable, colback=yellow!10, colframe=orange!70!black,
  title={Instructor note: what to draw}, fonttitle=\bfseries}

\title{Integration by Parts (Section 7.2) \\ \large Module Breakdown for AS.110.106 Calculus I (Biology and Social Sciences)}
\author{Concept-level teaching package \\ \small Three sessions, Module F (Week 12)}
\date{\today}

\begin{document}
\maketitle
\tableofcontents

%=================================================================
\section*{Assumptions}
\addcontentsline{toc}{section}{Assumptions}
\begin{itemize}
\item \textbf{Scope.} The module is Section 7.2 (Integration by Parts and Practicing Integration) of the syllabus, which sits inside Module F (Week 12). Sections 7.1 (substitution) and 7.3 (partial fractions) are not broken down here; 7.1 is treated as prior knowledge.
\item \textbf{Sessions.} Three sessions, all devoted to Section 7.2. The syllabus assumes only about 2--3 sessions for all of Week 12, so this plan implicitly assumes that 7.1 and 7.3 are scheduled around these three sessions. If all of Module F must fit into one week, Topic 3 (tabular and cyclic extensions) and Topic 6 are the natural places to compress.
\item \textbf{Level and prior knowledge.} Taken from the syllabus: introductory undergraduate calculus for biology and social-science majors. Students arrive having completed Modules A--E (through the Fundamental Theorem of Calculus and applications of integration) and Section 7.1.
\item \textbf{Topic emphasis.} None was specified, so the question bank is spread evenly across all six topics.
\item \textbf{Heuristic.} The syllabus lists ``LIATE/ILATE.'' This package uses LIATE; ILATE swaps the first two letters (inverse trigonometric before logarithmic), a difference that rarely matters at this level.
\item \textbf{Items beyond the syllabus.} Cyclic integrals (e.g.\ $\int e^x\sin x\,dx$) and inverse-trigonometric integrals (e.g.\ $\int \arctan x\,dx$) are included even though the syllabus does not list them by name. The cyclic case is marked optional. Inverse-trigonometric derivatives appear in Section 4.10, so the arctangent items are within reach; drop them if your course skips inverse-trigonometric integrals.
\item \textbf{Problems.} All examples are original. I did not have the textbook's own Section 7.2 exercises, so check notation against Neuhauser \& Roper. Every antiderivative and definite-integral value in this document was verified by differentiation or symbolic computation.
\item \textbf{Visual aids.} Described in words only, as instructor notes for sketching on the board. No images or diagram code are included.
\end{itemize}

%=================================================================
\section{Module Overview and Session Plan}

The module is broken into six topics, taught in the order below across three sessions.

\begin{longtable}{p{1.8cm}p{5.9cm}p{7.2cm}}
\toprule
\textbf{Session} & \textbf{Topics} & \textbf{Focus} \\
\midrule
\endfirsthead
\toprule
\textbf{Session} & \textbf{Topics} & \textbf{Focus} \\
\midrule
\endhead
\bottomrule
\endfoot
1 & Topic 1: The integration-by-parts formula \newline Topic 2: Choosing $u$ and $dv$ & Derive the formula from the product rule, evaluate first integrals, and show why the choice of $u$ decides whether the method helps. \\
2 & Topic 3: Repeated integration by parts and the tabular method \newline Topic 4: Logarithms and inverse trigonometric functions by parts & Extend the method to several passes, then to integrands with a single factor that cannot be integrated directly. \\
3 & Topic 5: Definite integrals by parts and applications \newline Topic 6: Choosing among integration techniques & Connect the method to the FTC and to accumulated change, then practice deciding among direct integration, substitution, and parts. \\
\end{longtable}

%=================================================================
\section{Prerequisites and Connecting to Prior Learning}

Prerequisites for each topic, and where each was taught:

\begin{longtable}{p{3.2cm}p{6.9cm}p{4.7cm}}
\toprule
\textbf{Topic} & \textbf{Prerequisite(s)} & \textbf{Where taught} \\
\midrule
\endfirsthead
\toprule
\textbf{Topic} & \textbf{Prerequisite(s)} & \textbf{Where taught} \\
\midrule
\endhead
\bottomrule
\endfoot
1. The formula & Product rule; derivatives of $e^x$ and trigonometric functions; antiderivatives of elementary functions; differential notation $dy=f'(x)\,dx$; checking an antiderivative by differentiating & 4.4; 4.8--4.9; 5.10; 4.11 \\
2. Choosing $u$ and $dv$ & Topic 1; power rule; derivative of $\ln x$; $\int e^{ax}\,dx$ and $\int\sin(ax)\,dx$ (undoing the chain rule) & This module (Topic 1); 4.3; 4.10; 5.10 and 7.1 \\
3. Repeated parts and tabular method & Topics 1--2; higher derivatives of polynomials; careful sign bookkeeping with parentheses & This module; 4.7; Ch.~1 (algebra background) \\
4. $\ln x$ and inverse trig & Topics 1--2; derivatives of $\ln x$ and $\arctan x$; substitution for $\int \frac{x}{1+x^2}\,dx$; logarithm rules & This module; 4.10; 7.1; Ch.~1 and 2.1 \\
5. Definite integrals & Topics 1--4; FTC Part 2; properties and signed-area meaning of the definite integral; accumulated change; the limit-conversion step of substitution (for contrast) & This module; 6.1--6.3; 7.1 \\
6. Choosing a technique & Substitution; the antiderivative table; spotting an inner function; Topics 1--5 & 7.1; 5.10; 4.5; this module \\
\end{longtable}

\begin{bridgebox}
\textbf{Where students are coming from.} They have just learned substitution, which the syllabus describes as the chain rule run in reverse. Integration by parts is the sequel: \emph{the product rule run in reverse}. Say this in the first minute so the module has a home in what they already know.

\textbf{Opening move for Session 1 (about five minutes, students do the work).}
\begin{enumerate}
\item Ask them to differentiate $xe^x$. This is a Section 4.4 product-rule problem; they get $e^x + xe^x$.
\item Ask for an antiderivative of $xe^x + e^x$. It is immediate: $xe^x$.
\item Ask for an antiderivative of $xe^x$ alone. Guide them to $xe^x = (xe^x+e^x) - e^x$, so $\int xe^x\,dx = xe^x - e^x + C$.
\item Name what they just did: they ran the product rule backwards and moved one piece to the other side. That is integration by parts; Topic 1 makes it general.
\end{enumerate}

\begin{tabular}{p{3.4cm}p{4.0cm}p{6.0cm}}
\multicolumn{3}{l}{\textbf{Two reversals, side by side (put this on the board).}} \\[4pt]
\toprule
\textbf{Rule (Module C)} & \textbf{Reversed technique} & \textbf{What to spot in the integrand} \\
\midrule
Chain rule (4.5) & Substitution (7.1) & An inner function whose derivative also appears as a factor. \\
Product rule (4.4) & Integration by parts & A product of two unrelated factors, one of which gets simpler when differentiated. \\
\bottomrule
\end{tabular}

\medskip
\textbf{Topic-specific callbacks.}
\begin{itemize}
\item \textbf{Topic 1 (Sections 4.4 and 4.11).} In 4.4 students learned that $(fg)'\neq f'g'$. The integral version of the same warning is $\int fg\,dx \neq \int f\,dx\int g\,dx$; Example Question 3 in Topic 1 uses the product rule to catch exactly this error. The differentials from 4.11 ($dy=f'(x)\,dx$) are the same objects as $du=u'(x)\,dx$ and $dv=v'(x)\,dx$ here, so remind them the notation is not new.
\item \textbf{Topic 2 (Sections 4.3, 4.10, and 7.1).} Ask ``what happens to $x$ when we differentiate it? when we integrate it?'' (it becomes $1$; it becomes $x^2/2$) and ``what happens to $\ln x$ when we differentiate it?'' (the logarithm disappears: $1/x$). That is the whole reason logarithms come first in LIATE. Connect the ``spot the inside function'' drills from 7.1 to ``spot the piece to differentiate'' drills here.
\item \textbf{Topic 3 (Section 4.7).} Differentiating a polynomial repeatedly until it reaches $0$ is Section 4.7 in disguise; the left column of the tabular method is exactly that list.
\item \textbf{Topic 4 (Sections 4.10 and 5.10).} The antiderivative table from 5.10 has no entry for $\ln x$. Integration by parts is the first technique that fills that gap. The derivatives of $\ln x$ and $\arctan x$ from 4.10 are precisely what make the method work, and the leftover integral $\int\frac{x}{1+x^2}\,dx$ is a 7.1 substitution.
\item \textbf{Topic 5 (Sections 6.2, 6.3, and 7.1).} The boundary term $[uv]_a^b$ is FTC Part 2 applied to $(uv)'$, so the definite formula is not a new fact. Accumulated change from 6.3 supplies the biological applications. Contrast explicitly with 7.1: in substitution the limits are converted; in parts they are not.
\item \textbf{Topic 6 (Sections 7.1 and 5.10).} Put $\int xe^{x^2}\,dx$ (substitution) beside $\int xe^{x}\,dx$ (parts) so students see that the same visible ingredients can call for different techniques. Point forward to 7.3: partial fractions is the next entry in the same toolkit, and the ``ask the decision questions again after every step'' habit carries over.
\end{itemize}
\end{bridgebox}

%=================================================================
\section{Topic-by-Topic Breakdown}

%-----------------------------------------------------------------
\subsection{Topic 1: The Integration-by-Parts Formula}

\textbf{Concepts}
\begin{itemize}
\item The product rule in differential form: $d(uv) = u\,dv + v\,du$
\item Integrating both sides to obtain $\int u\,dv = uv - \int v\,du$
\item Integration by parts as a \emph{trade}: one integral is exchanged for another, and the goal is a simpler one
\item The four ingredients: $u$, $du = u'(x)\,dx$, $dv = v'(x)\,dx$, and $v$
\item Constant of integration: $v$ needs no $+C$; a single $+C$ is added at the end
\item Verifying an antiderivative by differentiating it with the product rule
\item Why there is no ``product rule for integrals'': $\int fg\,dx \neq \int f\,dx \int g\,dx$
\end{itemize}

\begin{defbox}
\begin{itemize}
\item \textbf{Integration by parts formula.} ``If $u$ and $v$ are differentiable functions of $x$, then
\[ \int u(x)\,v'(x)\,dx \;=\; u(x)\,v(x) \;-\; \int v(x)\,u'(x)\,dx. \]
In differential notation, with $du = u'(x)\,dx$ and $dv = v'(x)\,dx$, this reads $\int u\,dv = uv - \int v\,du$.''
\item \textbf{The four pieces.} ``$u$ is the factor we will differentiate. $dv$ is everything else, including the $dx$, and it is the factor we will integrate. $du$ is $u'\,dx$. $v$ is one antiderivative of $dv$.''
\item \textbf{What the formula does.} ``Integration by parts does not evaluate an integral. It converts $\int u\,dv$ into $\int v\,du$. It helps exactly when the second integral is easier than the first.''
\item \textbf{No constant in $v$.} ``When we find $v$ from $dv$, we take any single antiderivative and do not add a constant. If we did, the extra constant would cancel in $uv - \int v\,du$. We add $+C$ once, at the very end.''
\end{itemize}
\end{defbox}

\textbf{Learning objectives} (students will be able to)
\begin{enumerate}
\item Derive the integration-by-parts formula from the product rule.
\item State the formula in both function and differential notation, and identify $u$, $du$, $dv$, and $v$ in a given application.
\item Apply the formula to evaluate integrals such as $\int xe^x\,dx$ and $\int x\cos x\,dx$.
\item Evaluate a proposed antiderivative by differentiating it with the product rule.
\item Explain, with a counterexample, why the integral of a product is not the product of the integrals.
\end{enumerate}

\textbf{Difficult concepts}
\begin{itemize}
\item \textbf{Formula errors.} The typical slips are specific: dropping the minus sign, forgetting that a second integral $\int v\,du$ remains, differentiating $v$ instead of integrating $dv$, and writing $u$ where $v$ belongs in the leftover integral. \emph{Mitigation:} derive the formula from the product rule in class rather than presenting it, require a four-line template ($u=$, $du=$, $dv=$, $v=$) before assembling the answer, and have students check every answer by differentiating.
\item \textbf{The product-of-integrals misconception.} Students who integrate each factor separately are transferring linearity to a product. \emph{Mitigation:} Example Question 3 below; differentiate the wrong answer with the product rule so the extra term is visible.
\end{itemize}

\textbf{Example questions with solutions}
\begin{enumerate}
\item \emph{(Derive)} Starting from the product rule, derive the integration-by-parts formula.

\textit{Solution.} The product rule states $\dfrac{d}{dx}\big[u(x)v(x)\big] = u'(x)v(x) + u(x)v'(x)$. Integrate both sides with respect to $x$. The left side integrates back to $uv$ (any constant is absorbed into the final $+C$):
\[ uv = \int u'v\,dx + \int uv'\,dx. \]
Solving for the second integral gives $\int uv'\,dx = uv - \int u'v\,dx$. Writing $du = u'\,dx$ and $dv = v'\,dx$, this is $\int u\,dv = uv - \int v\,du$.

\item \emph{(Recall)} Evaluate $\displaystyle\int xe^x\,dx$.

\textit{Solution.} Take $u = x$ and $dv = e^x\,dx$, so $du = dx$ and $v = e^x$. Then
\[ \int xe^x\,dx = xe^x - \int e^x\,dx = xe^x - e^x + C. \]
Check: $\dfrac{d}{dx}(xe^x - e^x) = e^x + xe^x - e^x = xe^x$.

\item \emph{(Analysis)} A student integrates each factor separately: $\int xe^x\,dx = \dfrac{x^2}{2}\cdot e^x + C$. Use the product rule to show that this is wrong.

\textit{Solution.} Differentiate the claimed answer: $\dfrac{d}{dx}\left[\dfrac{x^2}{2}e^x\right] = xe^x + \dfrac{x^2}{2}e^x$. This is not $xe^x$ (the extra term $\frac{x^2}{2}e^x$ does not vanish for $x\neq 0$). In general, if $F' = f$ and $G' = g$, then $(FG)' = fG + Fg$, which is not $fg$. So $\int fg\,dx \neq \int f\,dx\int g\,dx$; the product rule shows the correct relationship.
\end{enumerate}

\begin{notebox}
\textbf{Picture 1: the growing rectangle (Session 1, right after deriving the formula).} Draw a rectangle with the lower-left corner at the origin, horizontal side labeled $u$ and vertical side labeled $v$; write $uv$ inside it as its area. Extend the rectangle by a thin strip on the right of width $du$ and a thin strip along the top of height $dv$. Shade the right strip and label its area $v\,du$; shade the top strip and label its area $u\,dv$; draw the tiny corner square and label it ``$du\,dv$, negligible.'' Caption: ``the change in area is $u\,dv + v\,du$,'' which is the product rule in differential form.

\textbf{Picture 2 (optional, for the curious).} Draw axes with $v$ horizontal and $u$ vertical. Draw an increasing curve from the origin to a point $(V, U)$. Draw the rectangle with corners $(0,0)$, $(V,0)$, $(V,U)$, $(0,U)$. Shade the region below the curve and label it $\int u\,dv$; shade the region above the curve inside the rectangle and label it $\int v\,du$. Caption: ``the two areas fill the rectangle, so $\int u\,dv + \int v\,du = UV$'' (this is the formula when $u$ and $v$ both start at $0$).
\end{notebox}

%-----------------------------------------------------------------
\subsection{Topic 2: Choosing \texorpdfstring{$u$ and $dv$}{u and dv}}

\textbf{Concepts}
\begin{itemize}
\item Two requirements for a good choice: we must be able to integrate $dv$, and $\int v\,du$ must be simpler than $\int u\,dv$
\item Which functions get simpler when differentiated (polynomials, logarithms, inverse trigonometric functions) and which do not get worse when integrated ($e^{ax}$, $\sin ax$, $\cos ax$)
\item LIATE as a heuristic: what it says, why it usually works, and that it is a rule of thumb rather than a theorem
\item Previewing $du$ and $v$ for both candidate choices before committing
\item Computing $v$ sometimes requires undoing a chain rule (for example $\int\sin 3x\,dx$)
\item The same integral with a bad choice becomes harder, not easier
\end{itemize}

\begin{defbox}
\begin{itemize}
\item \textbf{A good choice of $u$ and $dv$.} ``A choice is good if two things hold: first, we can integrate $dv$; second, the new integral $\int v\,du$ is simpler than the one we started with.''
\item \textbf{LIATE.} ``LIATE orders function types: \textbf{L}ogarithmic, \textbf{I}nverse trigonometric, \textbf{A}lgebraic (powers and polynomials), \textbf{T}rigonometric, \textbf{E}xponential. As a heuristic, let $u$ be the factor whose type comes earliest in the list, and let $dv$ be the rest. It is a rule of thumb, not a theorem, and we will meet an integral where it fails.''
\end{itemize}
\end{defbox}

\textbf{Learning objectives}
\begin{enumerate}
\item Explain the two criteria that make a choice of $u$ and $dv$ good.
\item Apply LIATE to choose $u$ and $dv$ for products of two function types.
\item Compare two choices for the same integral and evaluate which one leads to a simpler integral.
\item Justify, in terms of what differentiation and integration do to each function type, why LIATE usually works.
\end{enumerate}

\textbf{Difficult concepts}
\begin{itemize}
\item \textbf{A poor choice produces a harder integral, and LIATE applied blindly hides why.} Students who cannot say \emph{why} $u=x$ beats $u=e^x$ cannot adapt when the heuristic fails. \emph{Mitigation:} the side-by-side comparison in Example Question 1, the ``would I rather differentiate this or integrate this?'' question, and the two-column chart in the instructor note below.
\item \textbf{Hidden chain rule in $v$.} For $dv=\sin(3x)\,dx$, students write $v=-\cos(3x)$ and lose the factor $\tfrac13$. \emph{Mitigation:} require students to differentiate their $v$ mentally before continuing; connect it to the 7.1 substitution $w=3x$.
\end{itemize}

\textbf{Example questions with solutions}
\begin{enumerate}
\item \emph{(Analysis)} For $\int xe^x\,dx$, a student sets $u = e^x$ and $dv = x\,dx$. Carry out one step and explain why this choice is unproductive.

\textit{Solution.} Then $du = e^x\,dx$ and $v = \dfrac{x^2}{2}$, so
\[ \int xe^x\,dx = \frac{x^2}{2}e^x - \int \frac{x^2}{2}e^x\,dx. \]
The new integral has a higher power of $x$ than the original, so it is harder. Differentiating $x$ lowers its power ($x\to 1$) while integrating it raises the power ($x\to x^2/2$); $e^x$ is unchanged by either. So $x$ belongs in $u$ and $e^x\,dx$ in $dv$, as in Topic 1.

\item \emph{(Application)} Evaluate $\displaystyle\int x\ln x\,dx$ and explain why $u=x$, $dv=\ln x\,dx$ is not an available choice at this point in the course.

\textit{Solution.} Take $u = \ln x$ and $dv = x\,dx$, so $du = \dfrac{1}{x}\,dx$ and $v = \dfrac{x^2}{2}$:
\[ \int x\ln x\,dx = \frac{x^2}{2}\ln x - \int \frac{x^2}{2}\cdot\frac1x\,dx = \frac{x^2}{2}\ln x - \int \frac{x}{2}\,dx = \frac{x^2}{2}\ln x - \frac{x^2}{4} + C. \]
The other choice fails the first criterion: it would require $v=\int\ln x\,dx$, which we cannot yet compute. It also makes the point that differentiating $\ln x$ removes the logarithm.

\item \emph{(Application)} Evaluate $\displaystyle\int x\sin(3x)\,dx$.

\textit{Solution.} Take $u = x$ and $dv = \sin(3x)\,dx$, so $du = dx$ and $v = -\tfrac13\cos(3x)$ (undo the chain rule; this is the 7.1 substitution $w=3x$). Then
\[ \int x\sin(3x)\,dx = -\frac{x}{3}\cos(3x) + \frac13\int\cos(3x)\,dx = -\frac{x}{3}\cos(3x) + \frac19\sin(3x) + C. \]
Check: the derivative of $-\frac{x}{3}\cos 3x$ is $-\frac13\cos 3x + x\sin 3x$, and the derivative of $\frac19\sin 3x$ is $\frac13\cos 3x$; the sum is $x\sin 3x$.
\end{enumerate}

\begin{notebox}
\textbf{Chart 1: the same integral, two choices (for Example Question 1).} Draw a table with two columns headed ``Option A: $u=x$, $dv=e^x\,dx$'' and ``Option B: $u=e^x$, $dv=x\,dx$.'' Use five rows labeled $du$, $v$, ``new integral,'' ``simpler than the original?'', and ``verdict.'' Fill in Option A as $dx$, $e^x$, $\int e^x\,dx$, ``yes,'' ``good''; fill in Option B as $e^x\,dx$, $x^2/2$, $\int\frac{x^2}{2}e^x\,dx$, ``no, higher power,'' ``bad.'' Have students predict the last two rows before you write them.

\textbf{Chart 2: what each LIATE type does under $d/dx$ and $\int$.} Draw a table with columns ``function type,'' ``after differentiating,'' ``after integrating,'' and ``goes in.'' Rows: $\ln x$ (becomes $1/x$, simpler; becomes $x\ln x - x$, more complex; $u$); $\arctan x$ (becomes $\frac{1}{1+x^2}$, algebraic; becomes $x\arctan x-\frac12\ln(1+x^2)$; $u$); $x^n$ (power drops; power rises; $u$ usually); $\sin x$ (same complexity; same complexity; $dv$); $e^x$ (unchanged; unchanged; $dv$). Fill the last column last and ask students to predict it. Caption: ``Put in $u$ what gets simpler when differentiated; put in $dv$ what is harmless to integrate.''
\end{notebox}

%-----------------------------------------------------------------
\subsection{Topic 3: Repeated Integration by Parts and the Tabular Method}

\textbf{Concepts}
\begin{itemize}
\item Repeated application: the leftover integral $\int v\,du$ is itself a product that needs parts again
\item For a polynomial of degree $n$ times $e^{ax}$, $\sin ax$, or $\cos ax$, each pass lowers the polynomial degree by one, so $n$ passes are needed
\item Bookkeeping: keep the leftover integral in brackets so the minus sign distributes correctly
\item The tabular method: a derivative column, an antiderivative column, diagonal products, and alternating signs
\item When the tabular method applies (a factor that differentiates to $0$) and when it does not (for example $x^2\ln x$)
\item \emph{Optional extension:} cyclic integrals, where the original integral reappears and is solved for algebraically
\end{itemize}

\begin{defbox}
\begin{itemize}
\item \textbf{Repeated integration by parts.} ``Repeated integration by parts means applying the formula again to the integral $\int v\,du$ when that integral is itself a product of the same kind. For $x^n$ times $e^{ax}$, $\sin ax$, or $\cos ax$, we need $n$ passes.''
\item \textbf{Tabular method.} ``The tabular method is a bookkeeping shortcut for repeated parts. Column D lists $u$ and its successive derivatives until we reach $0$. Column I lists $dv$ and its successive antiderivatives. Multiply each entry in D by the entry one row \emph{below} it in I, attach signs $+,-,+,-,\dots$, and add the products. Then add $+C$.''
\item \textbf{Cyclic integral (optional).} ``An integral is cyclic under parts if, after two applications, the original integral reappears on the right-hand side. We then treat it as an unknown and solve for it algebraically.''
\end{itemize}
\end{defbox}

\textbf{Learning objectives}
\begin{enumerate}
\item Apply integration by parts repeatedly to evaluate $\int x^n e^{ax}\,dx$ and $\int x^n\sin(ax)\,dx$ for $n=2,3$.
\item Apply the tabular method and explain why it reproduces the repeated-parts computation.
\item Evaluate whether the tabular method is appropriate for a given integrand.
\item \emph{(Optional)} Evaluate a cyclic integral such as $\int e^x\sin x\,dx$ by solving for the integral.
\end{enumerate}

\textbf{Difficult concepts}
\begin{itemize}
\item \textbf{Sign distribution.} After the first pass the answer has the form $x^2e^x - \int 2xe^x\,dx$, and the second pass produces a bracket that must be multiplied by $-1$. Students routinely lose the sign or the factor $2$. \emph{Mitigation:} keep the leftover integral in brackets until it is finished; use the tabular method as a cross-check.
\item \textbf{The tabular method as magic.} Students who use the table without understanding it apply it to $x^2\ln x$, where the derivative column never reaches $0$. \emph{Mitigation:} build the table for $\int x^2e^x\,dx$ side by side with the two-pass computation before introducing the shortcut, and state the condition ``one column must reach $0$.''
\item \textbf{Cyclic integrals look like failure.} When the original integral reappears, students think they are going in circles. \emph{Mitigation:} announce in advance that if the original integral comes back, that is good news, because it can be treated as an unknown $I$; check by differentiating.
\end{itemize}

\textbf{Example questions with solutions}
\begin{enumerate}
\item \emph{(Application)} Evaluate $\displaystyle\int x^2e^x\,dx$ using integration by parts twice.

\textit{Solution.} First pass: $u=x^2$, $dv=e^x\,dx$, so $du=2x\,dx$ and $v=e^x$:
\[ \int x^2e^x\,dx = x^2e^x - \int 2xe^x\,dx. \]
Second pass on $\int 2xe^x\,dx = 2\int xe^x\,dx = 2(xe^x - e^x)$ (Topic 1). Keeping the bracket:
\[ \int x^2e^x\,dx = x^2e^x - 2\big(xe^x - e^x\big) + C = x^2e^x - 2xe^x + 2e^x + C. \]
Check: the derivative is $2xe^x + x^2e^x - 2e^x - 2xe^x + 2e^x = x^2e^x$.

\item \emph{(Application)} Evaluate $\displaystyle\int x^3e^{2x}\,dx$ using the tabular method.

\textit{Solution.} Set up the table (signs are attached to the D entries; each D entry is multiplied by the I entry one row below it):

\smallskip
\begin{tabular}{cll}
\toprule
\textbf{Sign} & \textbf{D (differentiate)} & \textbf{I (integrate)} \\
\midrule
$+$ & $x^3$ & $e^{2x}$ \\
$-$ & $3x^2$ & $\tfrac12 e^{2x}$ \\
$+$ & $6x$ & $\tfrac14 e^{2x}$ \\
$-$ & $6$ & $\tfrac18 e^{2x}$ \\
 & $0$ & $\tfrac1{16} e^{2x}$ \\
\bottomrule
\end{tabular}

\smallskip
The diagonal products are $+x^3\cdot\tfrac12e^{2x}$, $-3x^2\cdot\tfrac14e^{2x}$, $+6x\cdot\tfrac18e^{2x}$, and $-6\cdot\tfrac1{16}e^{2x}$. Therefore
\[ \int x^3e^{2x}\,dx = e^{2x}\left(\frac{x^3}{2} - \frac{3x^2}{4} + \frac{3x}{4} - \frac38\right) + C. \]
Check: if $P(x)$ is the polynomial in parentheses, the derivative is $e^{2x}\big(2P+P'\big)$, and $2P + P' = x^3$.

\item \emph{(Optional extension: cyclic)} Evaluate $\displaystyle\int e^x\sin x\,dx$.

\textit{Solution.} Let $I=\int e^x\sin x\,dx$. Take $u=\sin x$, $dv=e^x\,dx$: $I = e^x\sin x - \int e^x\cos x\,dx$. On the new integral take $u=\cos x$, $dv=e^x\,dx$: $\int e^x\cos x\,dx = e^x\cos x + \int e^x\sin x\,dx = e^x\cos x + I$. Substituting,
\[ I = e^x\sin x - e^x\cos x - I \quad\Longrightarrow\quad 2I = e^x(\sin x - \cos x) \quad\Longrightarrow\quad I = \tfrac12e^x(\sin x - \cos x) + C. \]
(Keep the same type of factor in $u$ on both passes; switching roles on the second pass undoes the first and returns $I=I$.)
\end{enumerate}

\begin{notebox}
\textbf{Diagram: the tabular layout.} Draw three columns headed ``Sign,'' ``D,'' and ``I.'' Fill the D column from the top with $u$ and its successive derivatives, ending in $0$. Fill the I column from the top with $dv$ and its successive antiderivatives, one entry longer than D. Under Sign write $+,-,+,-,\dots$ next to the D entries. Draw a diagonal arrow from each D entry down and to the right to the I entry in the next row, and label each arrow with its sign. Stop when the D column hits $0$; circle the $0$ to show why the process ends. Next to it, write the two-pass calculation for $\int x^2e^x\,dx$ so students can see each diagonal product appear in it.
\end{notebox}

%-----------------------------------------------------------------
\subsection{Topic 4: Logarithms and Inverse Trigonometric Functions by Parts}

\textbf{Concepts}
\begin{itemize}
\item Integrands with a single factor that cannot be integrated directly: rewrite $\ln x$ as $1\cdot\ln x$ and take $dv=dx$, $v=x$
\item Why $\ln x$ and $\arctan x$ work as $u$: their derivatives, $\frac1x$ and $\frac1{1+x^2}$, are algebraic (Section 4.10)
\item Standard results: $\int\ln x\,dx$, $\int x^n\ln x\,dx$, $\int\arctan x\,dx$, $\int(\ln x)^2\,dx$
\item Computing $du$ correctly with the chain rule when the argument is not $x$ (for example $\ln(5x)$)
\item The leftover integral may need algebra or a 7.1 substitution to finish
\end{itemize}

\begin{defbox}
\begin{itemize}
\item \textbf{The $dv=dx$ move.} ``If the integrand is a single function $f(x)$ that we cannot integrate directly but can differentiate simply, such as $\ln x$ or $\arctan x$, write it as $1\cdot f(x)$. Take $u=f(x)$ and $dv=dx$. Then $v=x$.''
\item \textbf{Derivatives that make it work (from Section 4.10).} ``$\dfrac{d}{dx}\ln x = \dfrac1x$ and $\dfrac{d}{dx}\arctan x = \dfrac1{1+x^2}$. Differentiating turns each of these transcendental functions into a rational function, and that is why they are placed first in LIATE.''
\end{itemize}
\end{defbox}

\textbf{Learning objectives}
\begin{enumerate}
\item Explain why taking $dv=dx$ is legitimate for an integrand such as $\ln x$.
\item Evaluate $\int\ln x\,dx$, $\int x^n\ln x\,dx$, $\int\arctan x\,dx$, and $\int(\ln x)^2\,dx$.
\item Compute $du$ correctly when the argument of the logarithm or inverse trigonometric function is not simply $x$.
\item Analyze why the derivatives of $\ln x$ and $\arctan x$ turn the leftover integral into an algebraic one.
\end{enumerate}

\textbf{Difficult concepts}
\begin{itemize}
\item \textbf{``There is only one factor.''} Taking $dv=dx$ feels like cheating, because students expect two factors. \emph{Mitigation:} write $\int 1\cdot\ln x\,dx$ explicitly and ask ``what is an antiderivative of $1$?''
\item \textbf{The leftover integral needs more work.} For $\arctan x$ the second integral is $\int\frac{x}{1+x^2}\,dx$, which is not a parts problem; students conclude they made an error. \emph{Mitigation:} tell them in advance that the leftover integral may need substitution, and list it as an expected step.
\item \textbf{Chain rule inside $du$.} For $\ln(5x)$ students write $du=\frac{1}{5x}\,dx$ instead of $\frac1x\,dx$. \emph{Mitigation:} recall the 4.10 Recall exercise ($\frac{d}{dx}\ln(5x)=\frac1x$) and require $du$ to be computed on its own line.
\end{itemize}

\textbf{Example questions with solutions}
\begin{enumerate}
\item \emph{(Recall)} Evaluate $\displaystyle\int\ln x\,dx$.

\textit{Solution.} Write the integrand as $1\cdot\ln x$ and take $u=\ln x$, $dv=dx$, so $du=\frac1x\,dx$ and $v=x$:
\[ \int\ln x\,dx = x\ln x - \int x\cdot\frac1x\,dx = x\ln x - \int 1\,dx = x\ln x - x + C. \]

\item \emph{(Application)} Evaluate $\displaystyle\int\arctan x\,dx$.

\textit{Solution.} Take $u=\arctan x$, $dv=dx$, so $du=\dfrac{dx}{1+x^2}$ and $v=x$:
\[ \int\arctan x\,dx = x\arctan x - \int\frac{x}{1+x^2}\,dx. \]
For the leftover integral let $w=1+x^2$, $dw=2x\,dx$: $\int\frac{x}{1+x^2}\,dx=\frac12\int\frac{dw}{w}=\frac12\ln(1+x^2)$ (no absolute value is needed since $1+x^2>0$). Hence
\[ \int\arctan x\,dx = x\arctan x - \tfrac12\ln(1+x^2) + C. \]

\item \emph{(Application)} Evaluate $\displaystyle\int(\ln x)^2\,dx$.

\textit{Solution.} Take $u=(\ln x)^2$, $dv=dx$, so $du=\dfrac{2\ln x}{x}\,dx$ and $v=x$:
\[ \int(\ln x)^2\,dx = x(\ln x)^2 - \int x\cdot\frac{2\ln x}{x}\,dx = x(\ln x)^2 - 2\int\ln x\,dx. \]
Using Question 1, $\int\ln x\,dx = x\ln x - x$, so
\[ \int(\ln x)^2\,dx = x(\ln x)^2 - 2x\ln x + 2x + C. \]
\end{enumerate}

%-----------------------------------------------------------------
\subsection{Topic 5: Definite Integrals by Parts and Applications}

\textbf{Concepts}
\begin{itemize}
\item The definite formula $\int_a^b u\,dv = [uv]_a^b - \int_a^b v\,du$, obtained by applying FTC Part 2 to $(uv)'$
\item The limits stay in terms of $x$; unlike substitution, there is no limit conversion
\item The boundary term must be evaluated at \emph{both} endpoints, paying attention to the sign of the lower-limit contribution
\item Interpretation: signed area under a curve, and accumulated change from a rate such as $te^{-kt}$
\item Sanity checks: sign, magnitude, and a rough graph or estimate
\end{itemize}

\begin{defbox}
\begin{itemize}
\item \textbf{Definite integration by parts.} ``For differentiable $u$ and $v$ on $[a,b]$,
\[ \int_a^b u(x)\,v'(x)\,dx \;=\; \big[u(x)v(x)\big]_a^b \;-\; \int_a^b v(x)\,u'(x)\,dx, \]
where $\big[uv\big]_a^b = u(b)v(b) - u(a)v(a)$. The limits stay in terms of $x$; we do not change them.''
\item \textbf{Accumulated change (reminder from Section 6.3).} ``If $R(t)$ is a rate of change, then $\int_a^b R(t)\,dt$ is the total change accumulated from time $a$ to time $b$.''
\end{itemize}
\end{defbox}

\textbf{Learning objectives}
\begin{enumerate}
\item Derive the definite version of the formula from FTC Part 2 and the product rule.
\item Evaluate definite integrals by parts, keeping the limits in $x$ and evaluating $[uv]_a^b$ at both endpoints.
\item Interpret a definite integral computed by parts as a signed area or as accumulated change in a biological setting.
\item Evaluate the reasonableness of a computed value using a sketch or a rough numerical estimate.
\end{enumerate}

\textbf{Difficult concepts}
\begin{itemize}
\item \textbf{Two opposite errors with the limits.} Some students evaluate $uv$ only at the upper limit; others carry over the 7.1 habit and convert the limits to a new variable. \emph{Mitigation:} write $[uv]_a^b$ on the first line, evaluate the boundary term on its own line, and use a two-row contrast table (7.1: convert the limits; 7.2: keep the limits).
\item \textbf{Lower-limit sign.} In $\big[xe^x - e^x\big]_0^1$ the lower limit contributes $-1$, not $0$. \emph{Mitigation:} require each endpoint to be substituted in parentheses before simplifying.
\item \textbf{Trusting algebra over sense.} Students report a negative area or an impossible magnitude without noticing. \emph{Mitigation:} require a one-line reasonableness check against a sketch or a trapezoid estimate.
\end{itemize}

\textbf{Example questions with solutions}
\begin{enumerate}
\item \emph{(Recall)} Evaluate $\displaystyle\int_0^1 xe^x\,dx$.

\textit{Solution.} Take $u=x$, $dv=e^x\,dx$, so $du=dx$ and $v=e^x$ (limits stay in $x$):
\[ \int_0^1xe^x\,dx = \big[xe^x\big]_0^1 - \int_0^1e^x\,dx = (e - 0) - (e-1) = 1. \]

\item \emph{(Application)} Evaluate $\displaystyle\int_0^{\pi}x\sin x\,dx$ and interpret the value as an area.

\textit{Solution.} Take $u=x$, $dv=\sin x\,dx$, so $du=dx$ and $v=-\cos x$:
\[ \int_0^\pi x\sin x\,dx = \big[-x\cos x\big]_0^\pi + \int_0^\pi\cos x\,dx = \big(-\pi\cos\pi - 0\big) + \big[\sin x\big]_0^\pi = \pi + 0 = \pi. \]
Since $x\sin x\ge0$ on $[0,\pi]$, the value $\pi\approx3.14$ is the (unsigned) area of the region under the curve. It is plausible: the curve peaks at about $1.82$, so the bounding rectangle has area about $\pi\cdot1.82\approx5.7$, and the region fills a bit more than half of it.

\item \emph{(Application, biological)} A drug is absorbed into the bloodstream at rate $R(t)=te^{-t/2}$ mg/hr, $t$ hours after ingestion. How much is absorbed in the first 4 hours?

\textit{Solution.} The total absorbed is $\int_0^4te^{-t/2}\,dt$. Take $u=t$, $dv=e^{-t/2}\,dt$, so $du=dt$ and $v=-2e^{-t/2}$:
\begin{align*}
\int_0^4te^{-t/2}\,dt &= \big[-2te^{-t/2}\big]_0^4 + 2\int_0^4e^{-t/2}\,dt \\
&= -8e^{-2} + 2\big[-2e^{-t/2}\big]_0^4 = -8e^{-2} + 2\big(-2e^{-2}+2\big) = 4 - 12e^{-2}.
\end{align*}
So about $4-12(0.1353)\approx 2.38$ mg is absorbed. Reasonableness: on $[0,4]$ the rate stays between about $0.5$ and $0.74$ mg/hr (once it has left $0$), so a total near $4\times0.6\approx2.4$ mg fits.
\end{enumerate}

\begin{notebox}
\textbf{Graph A: $y=x\sin x$ on $[0,\pi]$.} Draw a horizontal axis labeled $x$ from $0$ to about $3.5$, with a tick at $\pi\approx3.14$; draw a vertical axis labeled $y$ from $0$ to $2$. Sketch the curve starting at the origin (it leaves the origin flat, like a parabola), rising to a maximum of about $1.82$ near $x\approx2.03$, then falling steeply to hit the axis at $(\pi,0)$. Shade the region under the curve and label it $\int_0^\pi x\sin x\,dx=\pi\approx3.14$. Lightly outline the bounding rectangle $[0,\pi]\times[0,1.82]$ to show the area is a bit over half of it.

\textbf{Graph B: the absorption rate $R(t)=te^{-t/2}$.} Draw a horizontal axis labeled ``$t$ (hours)'' from $0$ to $6$ and a vertical axis labeled ``$R$ (mg/hr)'' from $0$ to $0.8$. Plot the points $(0,0)$, $(1,0.61)$, $(2,0.74)$, $(3,0.67)$, $(4,0.54)$ and connect them with a smooth curve that peaks at $t=2$ and then decays slowly. Shade the region from $t=0$ to $t=4$ and label it ``total absorbed $=4-12e^{-2}\approx2.38$ mg.'' Emphasize that the shaded area, not the height of the curve, is the answer.

\textbf{Contrast table (7.1 versus 7.2).} Draw a two-row table: ``Substitution: limits are converted to the new variable; $x$ disappears.'' and ``Parts: limits stay in $x$; the boundary term $[uv]_a^b$ is evaluated at both ends.''
\end{notebox}

%-----------------------------------------------------------------
\subsection{Topic 6: Choosing Among Integration Techniques}

\textbf{Concepts}
\begin{itemize}
\item Decision questions asked in a fixed order: direct antiderivative, then substitution, then parts
\item Same ingredients, different technique: $\int xe^{x^2}\,dx$ (substitution) versus $\int xe^{x}\,dx$ (parts)
\item Combined problems: substitution first, then parts (or parts with a better choice of $u$)
\item LIATE is not always right: $\int x^3e^{x^2}\,dx$
\item Substitution can beat parts even when parts \emph{works}: $\int\frac{\ln x}{x}\,dx$ versus $\int\frac{\ln x}{x^2}\,dx$
\item Re-asking the decision questions after every step; recognizing when neither technique applies (forward pointer to 7.3)
\end{itemize}

\begin{defbox}
\begin{itemize}
\item \textbf{The decision questions.} ``Ask, in this order: (1) Can I integrate this directly from the antiderivative table? (2) Is one factor, up to a constant, the derivative of an inner function of the other? If so, use substitution. (3) Is the integrand a product of two unrelated factors, or a single function like $\ln x$ that I can differentiate simply, where one factor simplifies when differentiated and the other can be integrated? If so, use parts. After any step, ask the questions again about the new integral.''
\end{itemize}
\end{defbox}

\textbf{Learning objectives}
\begin{enumerate}
\item Classify integrands as direct, substitution, parts, or combined, and justify each classification in one sentence.
\item Design a solution path for an integral that requires both substitution and parts.
\item Evaluate a worked solution for wasted effort (for example, parts on $\int\frac{\ln x}{x}\,dx$) and propose the more efficient technique.
\item Solve a mixed set of integrals using appropriate techniques and verify each by differentiating.
\end{enumerate}

\textbf{Difficult concepts}
\begin{itemize}
\item \textbf{Defaulting to one technique.} Students use parts whenever they see a product, or substitution whenever they see a composition, and cannot tell $\int xe^{x^2}\,dx$ from $\int xe^{x}\,dx$. \emph{Mitigation:} classification-only drills with no computing, in the spirit of the 7.1 ``spot the inside function'' drills; the paired integrals in Example Question 1.
\item \textbf{Mixed problems: order and bookkeeping.} In $\int x^3e^{x^2}\,dx$, students convert $x^2$ to $w$ but leave stray $x$'s in the integrand. \emph{Mitigation:} require that the substitution be completed so that no $x$ remains (write $\frac12\int we^{w}\,dw$) before parts begins.
\item \textbf{LIATE as a law.} LIATE would say $u=x^3$ and $dv=e^{x^2}\,dx$, but $e^{x^2}$ cannot be integrated. \emph{Mitigation:} use this integral as the deliberate counterexample to the heuristic.
\end{itemize}

\textbf{Example questions with solutions}
\begin{enumerate}
\item \emph{(Classification)} Without computing, name the technique for each integral: (a) $\int xe^{x^2}\,dx$; (b) $\int xe^{x}\,dx$; (c) $\int x\cos(x^2)\,dx$; (d) $\int x\cos x\,dx$; (e) $\int x^3e^{x^2}\,dx$.

\textit{Solution.} (a) Substitution, $w=x^2$: the factor $x$ is (half) the derivative of the inner function $x^2$. (b) Parts, $u=x$: no inner function whose derivative appears. (c) Substitution, $w=x^2$, for the same reason as (a). (d) Parts, $u=x$. (e) Combined: substitute $w=x^2$ (only one $x$ is absorbed by $dw$; the other $x^2$ becomes $w$), leaving $\frac12\int we^w\,dw$, which needs parts.

\item \emph{(Application)} Evaluate $\displaystyle\int x^3e^{x^2}\,dx$.

\textit{Solution.} Let $w=x^2$, $dw=2x\,dx$, and write $x^3\,dx = x^2\cdot x\,dx = w\cdot\frac12\,dw$:
\[ \int x^3e^{x^2}\,dx = \frac12\int we^w\,dw. \]
Now parts on $\int we^w\,dw$ (Topic 1): $=we^w - e^w$. So
\[ \int x^3e^{x^2}\,dx = \tfrac12\big(we^w - e^w\big) + C = \tfrac12e^{x^2}(x^2-1) + C. \]
\emph{Alternative route.} Parts directly with $u=x^2$, $dv=xe^{x^2}\,dx$ (so $v=\frac12e^{x^2}$) gives $\frac12x^2e^{x^2} - \int xe^{x^2}\,dx = \frac12x^2e^{x^2} - \frac12e^{x^2} + C$, the same answer. Note that the LIATE choice $u=x^3$, $dv=e^{x^2}\,dx$ would fail, since $dv$ cannot be integrated.

\item \emph{(Analysis)} A student evaluates $\int\frac{\ln x}{x}\,dx$ by parts with $u=\ln x$ and $dv=\frac1x\,dx$ and reaches an equation containing the same integral again. Explain what happened, finish the computation, and give the faster method.

\textit{Solution.} With $u=\ln x$, $du=\frac1x\,dx$, $dv=\frac1x\,dx$, $v=\ln x$, the formula gives $I=(\ln x)^2 - \int\frac{\ln x}{x}\,dx = (\ln x)^2 - I$. The original integral came back (a cyclic integral), so solve for it: $2I=(\ln x)^2$, hence $I=\frac12(\ln x)^2+C$. The faster method is substitution: with $w=\ln x$, $dw=\frac1x\,dx$, the integral is $\int w\,dw=\frac{w^2}{2}=\frac12(\ln x)^2+C$. Here the factor $\frac1x$ is the derivative of the inner function $\ln x$, so decision question (2) applies before question (3).
\end{enumerate}

\begin{notebox}
\textbf{Flowchart: the decision questions.} Draw four boxes in a column, top to bottom, joined by downward arrows labeled ``no.'' Box 1: ``Can I integrate this directly (5.10 table)?'' with a sideways arrow labeled ``yes'' to a box ``done.'' Box 2: ``Is one factor (up to a constant) the derivative of an inner function of the other?'' with a sideways ``yes'' arrow to ``substitution (7.1).'' Box 3: ``Is it a product (or a lone $\ln x$ or $\arctan x$) where one piece simplifies when differentiated and the other can be integrated?'' with a sideways ``yes'' arrow to ``parts (choose $u$ with LIATE, then check with the two criteria).'' Box 4 (bottom): ``Other technique, such as partial fractions (7.3).'' Draw a long return arrow from the ``substitution,'' ``parts,'' and ``other'' boxes back up to the top of the column, labeled ``ask the questions again about the new integral.''
\end{notebox}

%=================================================================
\section{Module Question Bank}

\begin{enumerate}
\item \textbf{[Easy]} \emph{(Topic 1)} A student writes $\displaystyle\int x\cos x\,dx = x\sin x + \int\sin x\,dx = x\sin x - \cos x + C$. Find the error and give the correct answer.

\textit{Solution.} The formula is $uv - \int v\,du$, so the second term should be subtracted, not added. With $u=x$, $dv=\cos x\,dx$, $du=dx$, $v=\sin x$:
\[ \int x\cos x\,dx = x\sin x - \int\sin x\,dx = x\sin x + \cos x + C. \]
Differentiation exposes the student's answer: $\frac{d}{dx}(x\sin x-\cos x) = 2\sin x + x\cos x\neq x\cos x$, while $\frac{d}{dx}(x\sin x+\cos x)=x\cos x$.

\item \textbf{[Medium]} \emph{(Topic 2)} Evaluate $\displaystyle\int x^4\ln x\,dx$, and explain in terms of $du$ and $v$ why $u=\ln x$ is the better choice.

\textit{Solution.} Take $u=\ln x$, $dv=x^4\,dx$, so $du=\frac1x\,dx$ and $v=\frac{x^5}{5}$:
\[ \int x^4\ln x\,dx = \frac{x^5}{5}\ln x - \int\frac{x^5}{5}\cdot\frac1x\,dx = \frac{x^5}{5}\ln x - \frac{x^5}{25} + C. \]
Why: differentiating $\ln x$ removes the logarithm ($du=\frac1x\,dx$ is algebraic), so the leftover integral is a plain power. The alternative $u=x^4$, $dv=\ln x\,dx$ needs $v=x\ln x-x$ (itself a parts result), and the leftover integral $\int 4x^3(x\ln x-x)\,dx$ still contains a logarithm, so it is no simpler.

\item \textbf{[Medium]} \emph{(Topic 3)} Evaluate $\displaystyle\int x^2e^{-3x}\,dx$.

\textit{Solution.} Tabular method with D: $x^2,\,2x,\,2,\,0$ and I: $e^{-3x},\,-\frac13e^{-3x},\,\frac19e^{-3x},\,-\frac1{27}e^{-3x}$. The diagonal products with signs $+,-,+$ are
\[ x^2\Big(-\tfrac13e^{-3x}\Big) - 2x\Big(\tfrac19e^{-3x}\Big) + 2\Big(-\tfrac1{27}e^{-3x}\Big), \]
so
\[ \int x^2e^{-3x}\,dx = -e^{-3x}\left(\frac{x^2}{3} + \frac{2x}{9} + \frac{2}{27}\right) + C. \]
Check: with $P=\frac{x^2}{3}+\frac{2x}{9}+\frac2{27}$ the derivative is $e^{-3x}(3P-P')$, and $3P-P'=x^2$.

\item \textbf{[Medium]} \emph{(Topic 4)} Evaluate $\displaystyle\int\ln(5x)\,dx$.

\textit{Solution.} Take $u=\ln(5x)$ and $dv=dx$. By the chain rule $du=\frac{5}{5x}\,dx=\frac1x\,dx$, and $v=x$:
\[ \int\ln(5x)\,dx = x\ln(5x) - \int x\cdot\frac1x\,dx = x\ln(5x) - x + C. \]
Consistency check: $x\ln(5x)-x = x\ln x + (\ln5-1)x$, which agrees with integrating $\ln 5+\ln x$ term by term.

\item \textbf{[Medium]} \emph{(Topic 5)} Evaluate $\displaystyle\int_1^e x\ln x\,dx$.

\textit{Solution (rubric: correct choice of $u$; limits kept in $x$; both endpoints in the boundary term).} Take $u=\ln x$, $dv=x\,dx$, $du=\frac1x\,dx$, $v=\frac{x^2}{2}$:
\[ \int_1^ex\ln x\,dx = \Big[\tfrac{x^2}{2}\ln x\Big]_1^e - \int_1^e\frac{x}{2}\,dx = \tfrac{e^2}{2}(1) - 0 - \Big[\tfrac{x^2}{4}\Big]_1^e = \frac{e^2}{2} - \left(\frac{e^2}{4}-\frac14\right) = \frac{e^2+1}{4}\approx2.10. \]

\item \textbf{[Hard]} \emph{(Topic 5, graph-based)} A tracer is released into a stream at rate $r(t)=t^2e^{-t}$ mg/min, $t$ minutes after release. The graph of $r$ on $[0,3]$ is shown with the region under the curve shaded. (a) What does the shaded area represent? (b) Compute it exactly. (c) Use the graph to check that your value is reasonable.

\textit{Instructor note: add the image yourself.} Draw the graph of $r(t)=t^2e^{-t}$ for $0\le t\le3$: horizontal axis $t$ (minutes), vertical axis $r$ (mg/min) from $0$ to $0.6$; plot $(0,0)$, $(1,0.37)$, $(2,0.54)$ (the peak), $(3,0.45)$; connect with a smooth curve; shade the region between the curve and the $t$-axis from $0$ to $3$.

\textit{Solution.} (a) The total mass of tracer released during the first 3 minutes, in mg. (b) Take $u=t^2$, $dv=e^{-t}\,dt$, $du=2t\,dt$, $v=-e^{-t}$:
\[ \int_0^3t^2e^{-t}\,dt = \big[-t^2e^{-t}\big]_0^3 + 2\int_0^3te^{-t}\,dt. \]
For the second integral take $u=t$, $dv=e^{-t}\,dt$: $\int_0^3te^{-t}\,dt = \big[-te^{-t}\big]_0^3+\int_0^3e^{-t}\,dt = -3e^{-3} + \big[-e^{-t}\big]_0^3 = 1-4e^{-3}$. Therefore
\[ \int_0^3t^2e^{-t}\,dt = -9e^{-3} + 2\big(1-4e^{-3}\big) = 2 - 17e^{-3}\approx 1.15\text{ mg}. \]
(c) The curve stays below its peak of $0.54$ over a width of $3$, so the area is under $3\times0.54\approx1.62$; it fills roughly $70\%$ of that rectangle, and a trapezoid estimate from the four plotted points gives about $1.13$. Both agree with $1.15$.

\item \textbf{[Medium]} \emph{(Topic 6)} Evaluate each integral and name the technique: (a) $\displaystyle\int\frac{\ln x}{x}\,dx$; (b) $\displaystyle\int\frac{\ln x}{x^2}\,dx$.

\textit{Solution.} (a) Substitution: $w=\ln x$, $dw=\frac1x\,dx$, giving $\int w\,dw=\frac12(\ln x)^2+C$. (b) Parts: there is no inner function whose derivative appears, so take $u=\ln x$, $dv=x^{-2}\,dx$, $du=\frac1x\,dx$, $v=-\frac1x$:
\[ \int\frac{\ln x}{x^2}\,dx = -\frac{\ln x}{x} + \int\frac{1}{x}\cdot\frac{1}{x}\,dx = -\frac{\ln x}{x} - \frac1x + C. \]
The two integrals contain the same ingredients; only the presence or absence of a factor $\frac1x$ that is the derivative of $\ln x$ decides the technique.

\item \textbf{[Hard]} \emph{(Topic 6, combined with Topic 5)} Evaluate $\displaystyle\int_0^1x^5e^{x^2}\,dx$.

\textit{Solution (rubric: substitution completed before parts; limits converted once, at the substitution; two passes of parts done correctly).} Let $w=x^2$, $dw=2x\,dx$, so $x^5\,dx = x^4\cdot x\,dx = w^2\cdot\frac12\,dw$. The limits convert: $x=0\Rightarrow w=0$ and $x=1\Rightarrow w=1$. Then
\[ \int_0^1x^5e^{x^2}\,dx = \frac12\int_0^1w^2e^w\,dw. \]
Now parts twice (limits stay in $w$): $\int w^2e^w\,dw = w^2e^w - 2\int we^w\,dw = w^2e^w - 2\big(we^w-e^w\big) = e^w\big(w^2-2w+2\big)$. Evaluating,
\[ \frac12\Big[e^w\big(w^2-2w+2\big)\Big]_0^1 = \frac12\big(e\cdot1 - 1\cdot2\big) = \frac{e-2}{2}\approx0.36. \]
\end{enumerate}

\end{document}